% from ega4-2-fr.tex
\oldpage[IV-2]{47}

\begin{enumerate}
\item[$\alpha$)] $t$ est $M$-régulier~;
\item[$\beta$)] pour tout $i$ tel que $M_{\mathfrak p_i}\neq0$, l'image de $t$
dans $A/\mathfrak p_i$ engendre l'idéal maximal de cet anneau (ce qui entraîne
que $A/\mathfrak p_i$ est un anneau de valuation discrète).
\end{enumerate}

\begin{proof}
Comme $A$ n'est pas de dimension $0$ et que $A/tA$ est artinien, on a
nécessairement $\dim(A)=1$ $(0,16.3.4)$ et $t\notin\mathfrak p_i$ pour tout
$i$~: l'idéal principal $(t)$ est donc un idéal de définition de $A$, et
contient par suite une puissance de son idéal maximal $\mathfrak m$. Soit $N$
le sous-module des éléments de $M$ annulés par une puissance de $t$ (ou par
une puissance de $\mathfrak m$, ce qui revient au même comme on vient de le
voir)~; si l'on pose $P=M/N$, $t$ est $P$-régulier, car la relation $tx\in N$
pour un $x\in M$ entraîne $t^k(tx)=0$ pour un entier $k$, donc $x\in N$. Cela
étant, on a le

\begin{lemma}[3.4.1.4]
\label{IV.2.3.4.1.4-fr}
Soient $A$ un anneau,
\[
0\longrightarrow M'\longrightarrow M\longrightarrow M''\longrightarrow0
\]
une suite exacte de $A$-modules. Si $t\in A$ est $M''$-régulier, la suite
\[
0\longrightarrow M'/tM'\longrightarrow M/tM
\longrightarrow M''/tM''\longrightarrow0
\]
est exacte.
\end{lemma}

\begin{proof}
Comme $M/tM=M\otimes_A(A/tA)$, il suffit de prouver l'exactitude en $M'/tM'$~;
or, si l'image $x\in M$ d'un élément $x'$ de $M'$ est telle que $x=ty$ avec
$y\in M$, on en déduit, pour les images $x''$, $y''$ de $x$, $y$ dans $M''$,
$x''=ty''$~; mais comme $x''=0$, l'hypothèse entraîne $y''=0$, donc $y$ est
l'image d'un élément $y'\in M'$, et la relation $x=ty$ entraîne $x'=ty'$
puisque $M'\to M$ est injectif.
\end{proof}

Ce lemme étant établi, on en tire la relation
\end{proof}

\begin{env}[3.4.1.5]
\label{IV.2.3.4.1.5-fr}
\[
\operatorname{long}(M/tM)
=\operatorname{long}(N/tN)+\operatorname{long}(P/tP).
\]
\end{env}

D'autre part, pour tout $i$, on a $N_{\mathfrak p_i}=0$ puisque
$t\notin\mathfrak p_i$~; donc $M_{\mathfrak p_i}=P_{\mathfrak p_i}$~; pour
prouver (3.4.1.3), il suffit donc de le faire en y remplaçant $M$ par $P$~;
d'autre part, si les deux membres de (3.4.1.3) sont égaux, il résulte de la
même inégalité pour $P$ et de (3.4.1.5) que l'on a nécessairement
$\operatorname{long}(N/tN)=0$, donc $N/tN=0$ et finalement $N=0$, par le lemme
de Nakayama, $N$ étant de type fini~; or, $N=0$ signifie que $t$ est
$M$-régulier. On peut donc se ramener au cas où $M=P$, c'est-à-dire supposer
déjà que $t$ est $M$-régulier. Notons que cela entraîne que
$\mathfrak m\notin\operatorname{Ass}(M)$, puisque $t$ ne peut annuler un
élément $\neq0$ de $M$. Comme $A$ est de dimension $1$, on a donc
nécessairement
\[
\operatorname{Ass}(M)\subset\bigcup_i\{\mathfrak p_i\}.
\]

Procédons alors par récurrence sur
$n=\sum_i\operatorname{long}(M_{\mathfrak p_i})$. Si $n=0$, on a
nécessairement $M_{\mathfrak p_i}=0$ pour tout $i$, donc $M=0$ puisque aucun
des $\mathfrak p_i$ n'appartient à $\operatorname{Ass}(M)$~; les deux membres
de (3.4.1.3) sont alors nuls, et l'assertion $\beta$) de (3.4.1.2) est
triviale. Si $n>0$, le raisonnement du début de la démonstration de (3.4.1)
permet de supposer en outre que le $A$-module $M$ est fidèle~; cela entraîne
$M_{\mathfrak p_i}\neq0$ pour tout $i$ (Bourbaki, \emph{Alg. comm.}, chap.~II,
\S~2, n\textsuperscript{o}~2, cor.~2 de la prop.~4), et par suite
\[
\operatorname{Ass}(M)=\bigcup_i\{\mathfrak p_i\}.
\]

